% Part: first-order-logic
% Chapter: completeness
% Section: outline

\documentclass[../../../include/open-logic-section]{subfiles}

\begin{document}

\iftag{FOL}
      {\olfileid{fol}{com}{out}}
      {\olfileid{pl}{com}{out}}

\olsection{Garis Gedhe Bukti}

Bukti teorema kelengkapan rada rumit, lan nalika pisanan diwaca kita
gampang kesasar. Mula, ayo kita gawe garis gedhe buktine. Langkah
kapisan yaiku ngowahi sudut pandang supaya dalan tumuju bukti katon.
Nalika kelengkapan dirumusake minangka pratelan yen saben $\Gamma
\Entails !A$ mesthi menehi $\Gamma \Proves !A$, bisa uga angel
nemokake gagasan kanggo miwiti: kanggo nuduhake $\Gamma \Proves !A$,
kita kudu nemokake !!a{derivation}, dene hipotesis $\Gamma \Entails
!A$ katon ora maringi pitulungan kanggo nemokake iku. Ing sawetara
sistem bukti, !!a{derivation} bisa dibangun langsung, nanging kita
bakal nganggo ancangan kang rada beda. Owahing sudut pandang kang
dibutuhake yaiku iki: kelengkapan uga bisa dirumusake mangkene: yen
$\Gamma$ konsisten, mula himpunan iku bisa dicukupi. Mbokmenawa informasi
ing~$\Gamma$ bebarengan karo hipotesis konsistensine bisa digunakake
kanggo mbangun \iftag{FOL}{!!a{structure}}{!!a{valuation}} kang
nyukupi saben \iftag{FOL}{!!{sentence}}{!!{formula}} ing~$\Gamma$.
Kita wis ngerti jinis \iftag{FOL}{!!{structure}}{!!{valuation}} kang
digoleki: sing cocog persis karo katerangane~$\Gamma$!{}

Yen $\Gamma$ mung ngemot \iftag{FOL}{!!{sentence}s
  atomik}{!!{propositional variable}s}, modele gampang dibangun.
\iftag{FOL}{Upamana kabeh !!{sentence}s atomik mau awujud
  $\Atom{P}{a_1,\dots,a_n}$, kanthi $a_i$ minangka !!{constant}s.}{}
Kita mung perlu nemtokake
\iftag{FOL}{%
  !!a{domain}~$\Domain{M}$ lan pamasangan kanggo~$P$ saengga
  $\Sat{M}{\Atom{P}{a_1,\ldots,a_n}}$. Iki ora angel: tetepna
  $\Domain{M} = \Nat$, $\Assign{\Obj c_i}{M} = i$, lan kanggo saben
  $\Atom{P}{a_1,\ldots,a_n} \in \Gamma$, lebokna tuple $\tuple{k_1,
    \dots, k_n}$ menyang $\Assign{P}{M}$, kanthi $k_i$ minangka indeks
  simbol konstanta~$a_i$ (yaiku, $a_i \ident \Obj c_{k_i}$).}{%
  !!a{valuation}~$\pAssign{v}$ saengga $\pSat{v}{p}$ kanggo saben
  $p \in \Gamma$. Tetepna wae $\pAssign{v}(p) = \True$ yen lan mung
  yen $p \in \Gamma$.}

Saiki upamana $\Gamma$ ngemot !!{formula}~$\lnot !B$, kanthi $!B$
atomik. Bisa wae kita kuwatir yen konstruksi
$\iftag{FOL}{\Struct{M}}{\pAssign{v}}$ ngalangi $\lnot !B$ dadi
bener. Nanging ing kene konsistensi~$\Gamma$ nduweni peran: yen
$\lnot !B \in \Gamma$, mula $!B \notin \Gamma$; yen ora, $\Gamma$
ora konsisten. Yen $!B \notin \Gamma$, miturut konstruksi
$\iftag{FOL}{\Struct{M}}{\pAssign{v}}$ kita,
$\iftag{FOL}{\Sat/{M}}{\pSat/{v}}{!B}$, mula
$\iftag{FOL}{\Sat{M}}{\pSat{v}}{\lnot !B}$. Nganti kene ora ana
masalah.

Kepiye yen $\Gamma$ ngemot formula majemuk kang dudu atomik? Upamane,
himpunan iku ngemot $!A \land !B$. Supaya formula iku bener, kita kudu
tumindak kaya-kaya $!A$ lan $!B$ kalorone ana ing~$\Gamma$. Yen
$!A \lor !B \in \Gamma$, kita kudu ndadekake paling ora salah siji
ing antarane kalorone bener, yaiku tumindak kaya-kaya salah sijine ana
ing~$\Gamma$.

Iki nuwuhake gagasan mangkene: kita nambah !!{formula}s liyane
menyang~$\Gamma$ supaya (a)~himpunan asile tetep konsisten,
(b)~kanggo saben !!{sentence}~$!A$ ing basa kasebut, $!A$ utawa
$\lnot !A$ ana ing himpunan asil kasebut, lan (c)~yen $!A \land !B$
ana ing himpunan, $!A$ lan $!B$ kalorone uga ana; yen $!A \lor !B$
ana ing himpunan, paling ora salah siji saka $!A$ utawa $!B$ uga ana;
lan sateruse. Proses iki diterusake, bisa uga tanpa entek. Sebuten
himpunan kabeh !!{formula}s kang ditambahake mau~$\Gamma^*$. Banjur
konstruksi ing ndhuwur menehi
\iftag{FOL}{!!a{structure}~$\Struct{M}$}{!!a{valuation}~$\pAssign{v}$}
kang kanthi induksi bisa dibuktekake nyukupi kabeh ukara
ing~$\Gamma^*$, lan mula uga kabeh ukara ing~$\Gamma$ amarga
$\Gamma \subseteq \Gamma^*$. Cathetan koreksi sumber OLPL-045: sumber
Inggris nganggo wujud tunggal sentence nalika ngrujuk kabeh ukara ing
himpunan mau; wujud jamak kang dikarepake dibalekake ing kene. Pranyata
jaminan (a) lan~(b) wis cukup. Himpunan ukara kang nyukupi (b) diarani
\emph{jangkep}. Dadi, tugas kita yaiku ngambakake himpunan
konsisten~$\Gamma$ dadi himpunan~$\Gamma^*$ kang konsisten lan
jangkep. Cathetan koreksi sumber OLPL-737: syarat (b) nyakup kabeh ukara, manut definisi himpunan jangkep ing perangan sabanjure; watesan mung marang ukara atomik ing sumber ora cukup kanggo sipat kajangkepan kang diklaim.

\iftag{FOL}{%
Ana siji alangan ing rancangan iki: yen $\lexists[x][!A(x)] \in
\Gamma$, kita pengin bisa milih !!{constant}~$c$ lan nambah $!A(c)$
ing proses kasebut. Nanging kepiye kita ngerti manawa iki tansah bisa
dilakoni? Mbokmenawa basa kita mung duwe sawetara !!{constant}s, lan
kanggo saben konstanta mau $\lnot !A(c) \in \Gamma$. Kita ora bisa
uga nambah $!A(c)$, amarga iku bakal ndadekake himpunan ora konsisten,
lan kita ora bakal ngerti apa $\Struct{M}$ kudu ndadekake $!A(c)$
utawa $\lnot !A(c)$ bener. Kajaba iku, bisa wae $\Gamma$ mung ngemot
ukara ing basa kang ora duwe simbol konstanta babar pisan (umpamane
basa teori himpunan).

Cara ngrampungake masalah iki yaiku nambah tanpa winates akeh konstanta ing
wiwitan, bebarengan karo ukara-ukara kang nyambungake konstanta mau
marang kuantor kanthi cara kang trep. (Mesthi wae, kita kudu
mbuktekake yen iki ora bisa nuwuhake inkonsistensi.)

Konstruksi wiwitan kita lumaku kanthi becik yen ukara atomik mung
ngemot !!{constant}s. Nanging basa bisa uga ngemot !!{function}s.
Ing perkara iki, bisa angel nemokake fungsi ing~$\Nat$ kang trep
kanggo dipasangake marang !!{function}s kasebut. Mula ana cara liya:
tinimbang nganggo $i$ kanggo napsirake $\Obj c_i$, gunakna himpunan
!!{constant}s dhewe minangka domain. Banjur $\Struct M$ bisa
masangake saben !!{constant} marang awake dhewe:
$\Assign{\Obj c_i}{M} = \Obj c_i$. Malah, yagene ora diterusake
sakabehe: tetepna $\Domain{M}$ minangka kabeh \emph{term katutup} ing basa
kasebut!{} Cathetan koreksi sumber OLPL-738: domain nggunakake term katutup, cocog karo definisi model term ing perangan konstruksi model; watesan iki uga njaga supaya fungsi ing domain ngasilake term katutup. Kanthi mangkono, ana pamasangan fungsi kang cetha marang
!!{function}s: fungsi kasebut njupuk term minangka argumen lan
ngasilake term minangka nilai. Kita masangake marang
!!{function}~$\Obj f^n_i$ fungsi kang, yen diwenehi $n$ term
$t_1$, \dots,~$t_n$ minangka input, ngasilake term
$\Obj f^n_i(t_1, \dots, t_n)$ minangka nilai.

Perangan pungkasan yaiku carane nangani~$\eq$.
!!{predicate}~$\eq$ duwe interpretasi tetep:
$\Sat{M}{\eq[t][t']}$ yen lan mung yen $\Value{t}{M} =
\Value{t'}{M}$. Yen konstruksi digawe supaya !!{value} term~$t$
yaiku $t$ dhewe, !!{structure} iki ora bakal ndadekake
\emph{samubarang} ukara awujud $\eq[t][t']$ bener kajaba $t$ lan
$t'$ iku term kang padha persis. Iki dadi masalah, amarga meh saben
teori kang wigati ing basa kanthi !!{function}s duwe teorema awujud
$\eq[t][t']$ nalika $t$ lan $t'$ dudu term kang padha (umpamane ing
teori aritmetika: $\eq[(\Obj 0+ \Obj 0)][\Obj 0]$). Kanggo
ngrampungake masalah iki, domain~$\Struct M$ diowahi: tinimbang
nganggo term minangka objek ing~$\Domain{M}$, kita nganggo
himpunan-himpunan term, lan saben himpunan ngemot kabeh term kang
diwajibake padha dening ukara-ukara ing~$\Gamma$. Umpamane, yen
$\Gamma$ iku teori aritmetika, salah siji himpunan iki bakal ngemot
$\Obj 0$, $(\Obj 0 + \Obj 0)$, $(\Obj 0 \times \Obj 0)$, lan
sateruse. Himpunan iki bakal dipasangake marang $\Obj 0$, lan
pranyata uga dadi nilai kabeh term ing njero, kalebu
$(\Obj 0 + \Obj 0)$. Mula, ukara
$\eq[(\Obj 0+ \Obj 0)][\Obj 0]$ bakal bener ing
!!{structure} kang wis direvisi iki.}{}

Dadi, iki rancangan buktine. Kaping pisan, kita nliti sipat-sipat
himpunan konsisten kang !!{complete}; mligine, kita mbuktekake yen
himpunan konsisten kang !!a{complete} ngemot $!A \land !B$ yen lan mung
yen ngemot $!A$ lan~$!B$ kalorone, lan ngemot $!A \lor !B$ yen lan
mung yen ngemot paling ora salah sijine, lan sateruse
(\olref[ccs]{prop:ccs}).\iftag{FOL}{ Banjur kita nemtokake lan nliti
  himpunan ukara kang jenuh. Himpunan jenuh ngemot implikasi kang
  nyambungake saben !!{sentence} kang nganggo kuantor marang conto-conto
  ukara kasebut (\olref[hen]{defn:henkin-exp}). Kita mbuktekake yen
  saben himpunan konsisten~$\Gamma$ tansah bisa diambakake dadi
  himpunan jenuh~$\Gamma'$ (\olref[hen]{lem:henkin}). Yen sawijining
  himpunan iku konsisten, jenuh, lan !!{complete}, himpunan kasebut
  uga nduweni sipat yen ngemot
  \iftag{prvEx}{$\lexists[x][!A(x)]$ yen lan mung yen ngemot $!A(t)$
    kanggo sawijine term katutup~$t$}{}\iftag{defEx,defAll}{}{ lan
  }\iftag{prvAll}{$\lforall[x][!A(x)]$ yen lan mung yen ngemot $!A(t)$
    kanggo kabeh term katutup~$t$}{}
  (\olref[hen]{prop:saturated-instances}).}{}
Banjur kita njupuk himpunan\iftag{FOL}{ jenuh}{} kang
konsisten~$\iftag{FOL}{\Gamma'}{\Gamma}$ lan mbuktekake yen himpunan
iku bisa diambakake dadi himpunan kang
\iftag{FOL}{jenuh, konsisten, lan}{konsisten lan} !!{complete}~$\Gamma^*$
(\olref[lin]{lem:lindenbaum}). Himpunan~$\Gamma^*$ iki digunakake
kanggo nemtokake \iftag{FOL}{model term~$\Struct
  M(\Gamma^*)$}{!!{valuation}~$\pAssign v(\Gamma^*)$}.
\iftag{FOL}{Model term nduweni himpunan term katutup minangka domain,
  lan interpretasi !!{predicate}s ditemtokake dening !!{sentence}s
  atomik}{Valuasi ditemtokake dening !!{propositional
    variable}s}
ing~$\Gamma^*$ (\olref[mod]{defn:termmodel}). Kita nggunakake
sipat-sipat himpunan konsisten kang\iftag{FOL}{ jenuh,}{}
jangkep kanggo mbuktekake yen
$\iftag{FOL}{\Sat{M(\Gamma^*)}}{\pSat{v(\Gamma^*)}}{!A}$ yen lan mung
yen $!A \in \Gamma^*$ (\olref[mod]{lem:truth}), lan mligine,
$\iftag{FOL}{\Sat{M(\Gamma^*)}}{\pSat{v(\Gamma^*)}}{\Gamma}$.
\iftag{FOL}{Pungkasan, kita bakal nliti carane nemtokake model term
  yen $\Gamma$ uga ngemot~$\eq$
  (\olref[ide]{defn:term-model-factor}) lan mbuktekake yen model iku
  nyukupi~$\Gamma^*$ (\olref[ide]{lem:truth}).}{}

\end{document}
